\scanpage{32}
\mspage{15}
5. \uline{Spécialisation \deleted{et générisation} \unsure{de|des} MD-graphes.}

Soit
\figure{\begin{tikzpicture}[baseline=(f.base)]
\node[inner sep=7pt] (f) {$(38)\quad G=(S,\widetilde A,\sigma_{\widetilde A},\deleted{\unsure{\widetilde A|A}}I,\ \underset{\substack{\Vert\\\widetilde A\amalg I}}{\widehat I}\to S\xrightarrow[g]{}\deleted{\unsure{g}}\mathbb N)$};
\draw ([xshift=-.5cm]f.south west)--(f.south west)--(f.north west)--(f.north east)--(f.south east)--([xshift=.1cm]f.south east);
\end{tikzpicture}}
un MD-graphe, de type numérique $(g,\nu)$. Nous allons construire ses \deleted{$\unsure{G}$} spécialisations, comme les \inserted{MD-}graphes associés aux diverses ``strates'' de
\[
\deleted{\unsure{\widehat{\mathcal M}_G}}\qquad\widehat{\mathcal M}_G=\prod_{\alpha\in S}\widehat{\mathcal M}_{g_\alpha,\widehat I_\alpha}.
\]
Une telle strate est obtenue en choisissant, pour $\forall\ \alpha$, une strate de $\widehat{\mathcal M}_{g_\alpha,\widehat I_\alpha}$, c'est à dire une classe d'isomorphie de MD-graphe
\deleted{$G_\alpha=(S_{(\alpha)},\widetilde A_{(\alpha)},\sigma_{\widetilde A_{(\alpha)}},\ I_{(\alpha)}=\widehat I_\alpha\to S_{(\alpha)}\to\mathbb N)$ de type numérique}
de type \deleted{$(g_\alpha,\widehat\nu_\alpha)$} numérique $(g_\alpha,\widehat I_\alpha)$
\[
(39)\quad\deleted{\unsure{(g_\alpha,\widehat\nu_\alpha)}}.\quad
\begin{aligned}
G_{(\alpha)}=(S_{(\alpha)},\widetilde A_{(\alpha)},\sigma_{\widetilde A_{(\alpha)}},\ &I_{(\alpha)}=\widehat I_\alpha,\\
&\deleted{\unsure{\widehat I_{(\alpha)}\amalg\widetilde A_{(\alpha)}\to S_{(\alpha)}}\xrightarrow{\unsure{\sigma_\alpha|g_\alpha}}\mathbb N}\\
&\inserted{\widehat I_{(\alpha)}\overset{\mathrm{déf}}{=}\widetilde A_{(\alpha)}\amalg I_{(\alpha)}\to S_{(\alpha)}\xrightarrow{\unsure{\sigma_\alpha|g_\alpha}}\mathbb N})
\end{aligned}
\]
(\uline{NB} bien entendu, le cas où $G_{(\alpha)}$ est le MD-graphe trivial de type $(g_\alpha,\widehat\nu_\alpha)$ n'est pas exclu).

Par ``classes d'isomorphie'', il faut entendre ici : pour des isomorphismes induisant l'identité sur $I_{(\alpha)}=\widehat I_\alpha$. Au lieu de travailler avec de telles classes d'isomorphie, nous allons pour commencer travailler
\mspage{31}
